% =============================================================================
% 3 PRODUCT OWNER, INITIAL BACKLOG, REFINEMENT            budget: 2.50 pages
% Answers TASK QUESTION 2, verbatim:
%   "Put yourself in the role of the product owner and describe how the initial
%    backlog is created based on the defined user stories and how it is refined
%    with the Scrum team."
%
% THIS IS THE CORE CHAPTER. "This bit of the task probably will take a couple
% of pages" [G-Jan26]. The wording "refined with the Scrum team" is what
% distinguishes Task 2 from Task 1 - it must be visibly answered.
%
% WHAT TASK 2 ACTUALLY WANTS - his framing has settled:
%   "The difference between task one and task two is negligible... the only
%    difference is that the last sub question is, how is the product backlog
%    refined with the scrum team? Really, the only thing that you need to look
%    at here is describe how sprint planning happens" [G-Jan26]
%   "The main way in which this product backlog changes is through a process
%    called product backlog refinement... and mostly it's done through a
%    ceremony called sprint planning, or during an activity called backlog
%    refinement. So you need to research what these are and describe what these
%    are" [G-Jun26]
%   => COVER BOTH EVENTS: backlog refinement AND sprint planning, refinement
%      first. That also satisfies the older "refinement is the emphasis"
%      reading [G-Mar25, G-Nov25, S4].
%
% SECTION PLAN he gave himself [G-Dec25]: (a) the role as PO, (b) how
% requirements are gathered, (c) what the backlog is and how it changes,
% (d) prioritisation, (e) refinement - plus (f) sprint planning. The four
% subsections below cover all six.
%
% PERSPECTIVE: argue from the product owner's point of view - "it doesn't
% matter if you use the first person... but do think from the perspective of
% product owner" [G-Feb25]. Since [GL] forbids first person, write it in the
% third person while keeping the PO's viewpoint.
%
% THE #1 SCORING RULE APPLIES HERE MOST OF ALL: "don't just give a general
% description of what a product owner does, for example, but what he or she
% did in your specific project" [SCH, 01.10.23]. Every theoretical statement
% needs a recipe-app consequence in the same or the next sentence.
%
% TOOLS (Jira, Trello, Miro, Kanban board) may be mentioned IN TEXT ONLY, never
% as a figure. The reason to cite is transparency: "everything is very
% transparent in the world of agile development... everybody should know what's
% going on at any time... and then avoid miscommunication" [G-Feb25].
% =============================================================================

\section{Product Owner: Initiales Backlog und Refinement}

% =============================================================================
% 3.1 THE ROLE OF THE PRODUCT OWNER
% -----------------------------------------------------------------------------
% LITERATURE: HEAVY  ->  the second citation-dense passage. The PO is a defined
%                        role in a published framework, so every
%                        accountability claimed must be traceable.
%   Cite here:
%     - the Scrum Guide for the accountabilities (product goal, creating and
%       ordering PBIs, backlog transparency) - one page-numbered citation
%       covering the list is enough, do not cite it five times
%     - Kelly (2019) for the business-facing, negotiating side of the role;
%       this is the book the task list names for exactly this purpose
%     - the Agile Practice Guide for how the role sits between business and
%       team
%     - Schwaber's Scrum book, if the "no project manager in Scrum" sentence
%       needs backing (it already sits in literature/)
%   Research, best case: one source on the authority question - who may change
%   the backlog and why a single accountable person beats a committee. If none
%   is found, argue it from the Scrum Guide's wording instead of asserting it.
%   NOT cited: the "most important word is no" framing and the concrete
%   decisions your PO made in this project - those are your own voice.
%
% WHAT YOU WANT TO SAY HERE, IN SHORT
%   - "The PO is a single accountable person, the voice of the customer, and
%      these are the accountabilities the framework assigns." [cite]
%   - "It is a senior business role, not a technical one, and its core work is
%      negotiation, because demand always exceeds what the project can fund."
%      [cite Kelly]
%   - "The characteristic move is not refusal but pricing: here is the impact
%      if you want this in." [own voice]
%   - "Only the PO changes the backlog; anyone may propose changes." [cite]
%   - "There is no project manager in Scrum." [one sentence, cite]
%   - Throughout: what THIS PO decided for the recipe app, not what a PO does
%     in general.
% -----------------------------------------------------------------------------
% MATERIAL, in rough order of value:
%   - Voice of the customer [G-Mar25, G-Nov25].
%   - ONE PERSON, never a committee [S2].
%   - The Scrum Guide accountabilities as the backbone [S2, S4]: developing and
%     communicating the product goal, creating and clearly communicating PBIs,
%     ordering them, ensuring the backlog is transparent, visible and
%     understood. Cite the Scrum Guide with a page number.
%   - A senior business person with negotiating skills, often not a
%     technologist, who must make constraints plain because "there are going to
%     be many, many more requirements than the project can afford"
%     [G-Jan25, G-Feb25].
%   - THE KEY TRAIT, and the strongest single line available for the Resources
%     criterion: "the most important word in systems development is no" - with
%     the agile refinement "In fact, Agile doesn't say no. It says, well,
%     here's the impact, if you want this in" [G-Feb25].
%   - ONLY THE PO MAY CHANGE THE BACKLOG. Anyone may suggest additions, changes
%     or reprioritisation, "but only the product owner has the authority to
%     make the changes. So it isn't a random free for all" [G-Jan25].
%   - ONE SENTENCE: there is no project manager in Scrum - "There is project
%     management activity, but it's done on a kind of collective basis"
%     [G-Jun26]. Cheap Transfer point, do not expand it.
%
% SHAPE: do not let this become a role portrait. Two thirds role, one third
% what that role decided in THIS project.
% =============================================================================
% #############################################################################
% WRITING PLAN - 3.1, PARAGRAPH BY PARAGRAPH
% Budget: 0.60 pages = ~420 WORDS across 6 paragraphs => ~70 WORDS EACH.
% HEAVY - about 11 citations, but SHORT. That density is only affordable
% because most citations are a clause, not a block quote.
% Shape: two thirds what the role is, one third what THIS PO decided.
%
% Sources that own this subsection: scrumguide2020 (normative backbone),
% kelly2019 (the business and negotiating side), pmi2017 (the role between
% business and team), suryaatmaja2020 (what happens when the role is undefined).
%
% !! THREE SOURCES MUST BE KEPT OUT OF THIS SUBSECTION:
%    - cohn2004: he writes "A ScrumMaster fills the project manager role".
%    - schwaber2004: p. 1 calls the ScrumMaster "the Scrum project manager".
%    - racheva2008: its Table 2 has the "Scrum master (or Agile PM)" advising on
%      technical feasibility.
%    The literature review rejects two other sources for exactly this error, so
%    citing any of these next to "there is no project manager in Scrum" would be
%    inconsistent with the report's own screening.
% #############################################################################
\subsection{Die Rolle des Product Owners}

% --- P1: a single accountable person, and the four accountabilities -------
% THREE CITATIONS.
%
%   \parencite[p.~5]{scrumguide2020}
%   "The Product Owner is accountable for maximizing the value of the product
%    resulting from the work of the Scrum Team."
%
%   \parencite[p.~6]{scrumguide2020}
%   The four accountabilities, as a list: developing and explicitly
%   communicating the Product Goal; creating and clearly communicating Product
%   Backlog items; ordering Product Backlog items; ensuring the Product Backlog
%   is transparent, visible and understood.
%   -> ONE page-numbered citation covering the whole list. Do not cite p. 6 four
%      times. !! If quoted verbatim, strip the trailing "; and," from "Ordering
%      Product Backlog items; and,".
%   -> ** THE STRUCTURAL SPINE OF THIS SUBSECTION **: map each of the four onto
%      a concrete recipe-app action in the same or the next sentence. That
%      discharges the module's most repeated scoring rule ("not what a product
%      owner does in general, but what he or she did in your project").
%
%   \parencite[p.~6]{scrumguide2020}
%   "The Product Owner is one person, not a committee."
%   -> The strongest one-line quote available for 3.1.
%   -> A second, independent source exists if reinforcement is wanted:
%      \parencite[p.~153]{pmi2017} defines the PO as "A person responsible for
%      maximizing the value of the product and who is ultimately responsible and
%      accountable for the end product that is built." (Note: the literature
%      review's first pass wrongly listed PMI as silent on this.)
Der Product Owner ist \enquote{accountable for maximizing the value of the
product resulting from the work of the Scrum Team} \parencite[p.~5]{scrumguide2020}
und ist \enquote{one person, not a committee} \parencite[p.~6]{scrumguide2020}.
Daneben nennt der Guide eine zweite Rechenschaftspflicht, das wirksame Management
des Product Backlogs, das er in vier Teile gliedert: das Produktziel zu
entwickeln und zu vermitteln, Product-Backlog-Einträge zu erstellen und zu
vermitteln, sie zu ordnen und das Backlog transparent, sichtbar und verstanden zu
halten \parencite[p.~6]{scrumguide2020}. Der Product Owner darf diese Arbeit
delegieren, bleibt aber dafür rechenschaftspflichtig. Für Cookbox sind diese vier Punkte: das Ziel, dass
eine kochende Person ein Rezept in die App bekommt, es wiederfindet und ihm beim
Kochen folgt; die zehn Stories aus Tabelle~\ref{tab:user-stories}; ihre
Reihenfolge in Abschnitt~\ref{subsec:backlog}; und ein Backlog, das im gesamten
Studio einsehbar ist.

Diese Person gehört auf die Geschäftsseite. In der Regel gilt: Product Owner
\enquote{have a business background and bring deep subject matter expertise to
the decisions} \parencite[p.~41]{pmi2017}, und um wirksam zu sein, müssen sie
\enquote{step away from the keyboard and stop writing code} \parencite[ch.~11, pp.~79--84]{kelly2019};
der Product Owner von Cookbox ist eine leitende Geschäftsrolle im Studio und
schreibt keine Zeile der App. Accountability \enquote{is not shared}, während
responsibility \enquote{can be shared} \parencite[The Standard for Project Management,
Section~3.2]{pmbok2021}: Für die Reihenfolge der zehn Stories ist allein der
Product Owner rechenschaftspflichtig, während ihre Formulierung und ihre Größe
gemeinsam mit den Entwicklern erarbeitet werden.

Ein großer Teil dieser Arbeit ist Verhandlung: Sechs Monate und feste
240.000~EUR kaufen sehr viel weniger, als das Studio gern ausliefern würde.
\textcite[ch.~2, pp.~11--20]{kelly2019}
benennt diese Einschränkung unumwunden: \enquote{every work item which is done comes at
the cost of not doing something else.} Das wichtigste Wort dieser Rolle ist
deshalb \enquote{no}; agil wird jedoch nicht rundheraus abgelehnt, sondern
bepreist, indem die Auswirkung benannt wird, die eine Aufnahme hätte. Der Import
von Webseiten (US-10) wurde wiederholt eingefordert und wurde bepreist statt
abgelehnt, wobei der Preis am Releasetermin bemessen wurde. Eine einzige
ordnende Stimme hält das Backlog außerdem frei von \enquote{decibel management},
also jenem Zustand, in dem sich der lauteste Stakeholder durchsetzt
\parencite[ch.~2, pp.~11--20]{kelly2019}.

Vorschläge darf jeder machen, aber nur mit Argumenten: Wer das Product Backlog
verändern will, \enquote{can do so by trying to convince the Product Owner}
\parencite[p.~6]{scrumguide2020}. Diese Befugnis schließt das Streichen ein, denn
der Product Owner entscheidet auch \enquote{what should not be developed into product}
\parencite[ch.~8, pp.~57--62]{kelly2019}. Genau das hält das Won't-have-Band in
Tabelle~\ref{tab:initial-backlog} fest: US-10 steht dort als Entscheidung und
nicht als Versäumnis.

Innerhalb des Scrum Teams benennt das Framework nur drei Rechenschaftspflichten
--- Developers, Product Owner und Scrum Master --- \parencite[p.~5]{scrumguide2020}.
Ein Projektmanager gehört nicht dazu, und für Cookbox wurde auch keiner
eingesetzt; die Planung, die dieses Kapitel dokumentiert, haben diese drei
gemeinsam geleistet.

Der Zuschnitt der Rolle wurde deshalb vor dem ersten Sprint schriftlich
festgehalten: Der Product Owner ordnet das Backlog allein, nimmt an beiden
Refinement-Sitzungen jedes Sprints teil und verteidigt das Produktziel. Eine
aktionsforschende Studie zu einer gescheiterten Agil-Einführung hält fest, dass
\enquote{the team has learned that there were no clear
rules/guidelines for PO to follow} \parencite[p.~15]{suryaatmaja2020}. Eine
Rolle zu benennen und festzulegen, worüber sie entscheidet, sind zwei
verschiedene Handlungen, und dieses Projekt brauchte die zweite.


% --- P2: a business role, and the accountability/responsibility split -----
% THREE CITATIONS.
%
%   \parencite[p.~41]{pmi2017}
%   "Typically, product owners have a business background and bring deep subject
%    matter expertise to the decisions."
%   and, same page, "Product owners rank the work based on its business value."
%   -> Seven words that give both the act and the criterion. Use the second as
%      the hinge into 3.2.
%
%   \parencite[ch.~11, pp.~79--84]{kelly2019}
%   "But to be an effective Product Owner, they need to step away from the
%    keyboard and stop writing code."
%   -> The business-not-technologist point in the words of the book the task
%      sheet names for exactly this purpose.
%
%   \parencite[The Standard for Project Management, Section~3.2]{pmbok2021}
%   "Accountability. The condition of being answerable for an outcome.
%    Accountability is not shared." / "Responsibility. The condition of being
%    obligated to do or fulfill something. Responsibility can be shared."
%   -> ** USE THESE TWO TOGETHER **: the PO is ACCOUNTABLE for the product,
%      while the responsibility for refining the backlog is SHARED with the
%      Scrum team. That single distinction sets up 3.3 and is exactly the
%      reflection depth Process and Quality reward.
%   !! THE PART NAME IS MANDATORY IN THIS LOCATOR. The volume binds two
%      separately numbered documents and both have a Section 3.2; without the
%      prefix an examiner checking it lands on "Why Tailor?" and concludes the
%      quote is fabricated.


% --- P3: negotiation, and the characteristic move -------------------------
% TWO CITATIONS plus the report's own framing.
%
%   \parencite[ch.~2, pp.~11--20]{kelly2019}
%   "But every work item which is done comes at the cost of not doing something
%    else."
%   -> ** THE BEST AVAILABLE BACKING for the tutor's "here is the impact if you
%      want this in" framing. ** It states the PO's move as PRICING - every yes
%      is priced in foregone work - rather than as refusal. Preferred over
%      Kelly's "resources are limited... decide what is not being done"
%      (ch. 9, pp. 63-74), which frames it as refusal.
%
%   \parencite[ch.~2, pp.~11--20]{kelly2019}
%   "Having a 'single voice' creates clarity and forces prioritization. Many a
%    project has been waylaid by diverse stakeholders pulling the team in
%    different directions. Too often 'decibel management' has ruled."
%   -> "Decibel management" is the most quotable phrase in the book and the best
%      justification for why the backlog is ordered by one person rather than by
%      whoever shouts loudest. Also a Creativity point.
%
% OWN VOICE, UNCITED: the sentence that the most important word is "no", and
% its agile correction - agile does not refuse, it prices ("here is the impact
% if you want this in"). This is the tutor's own framing and must NOT be
% attributed to a source. It is the strongest single line for the Resources
% criterion (10 %).
% Then, concretely: name one request this PO priced rather than refused.


% --- P4: who may change the backlog ---------------------------------------
% TWO CITATIONS.
%
%   \parencite[p.~6]{scrumguide2020}
%   "Those wanting to change the Product Backlog can do so by trying to convince
%    the Product Owner."
%   -> Anyone may propose; only the PO decides. Contrast this in chapter 4 with
%      formal change control - do not do the contrast here.
%
%   \parencite[ch.~8, pp.~57--62]{kelly2019}
%   "At its most basic level, the Product Owner is the backlog gatekeeper... More
%    importantly the Product Owners should get to decide what is taken OUT of the
%    backlog, that is, what should not be developed into product."
%   -> The removal authority is the one thing the Scrum Guide does NOT state, so
%      this citation is not redundant. It also sets up the "Won't have" band in
%      3.2.


% --- P5: no project manager in Scrum (ONE sentence) -----------------------
% ONE CITATION. Do not expand this beyond a sentence - it is a cheap Transfer
% point that becomes expensive if it turns into a paragraph.
%
%   \parencite[p.~5]{scrumguide2020}
%   "Scrum defines three specific accountabilities within the Scrum Team: the
%    Developers, the Product Owner, and the Scrum Master."
%   -> Cite the list, then note the ABSENCE in your own words: a project-manager
%      role is not among them. Phrase it as an observation about the framework,
%      not as a quotation - the Guide never says "there is no project manager".
%      (Verified: "project manager" occurs 0 times in the Guide.)
%   -> Kelly's nearest sentence ("the Project Manager role steps back") is a
%      claim about EMPHASIS, not composition - do not use it for this.
% OPTIONAL half-sentence of balance, if space allows:
%   \parencite[p.~37]{pmi2017} notes that project managers "can add significant
%   value in many situations" but that their role looks different. That nuance
%   is a Quality point, but it costs a citation the budget may not have.


% --- P6: what THIS PO did, and the risk the role carries ------------------
% ONE CITATION. Close the subsection on the project, not on theory.
%
% OWN VOICE: two or three sentences on decisions this PO actually took for the
% recipe app - which stakeholder group was pushed back, which request was
% priced, what the product goal was.
%
%   \parencite[p.~15]{suryaatmaja2020}
%   "The team has learned that there were no clear rules/guidelines for PO to
%    follow"
%   -> ** WHY THIS SOURCE EARNS ITS PLACE HERE **: normative sources structurally
%      CANNOT report a failed adoption. This is documented evidence of what
%      happens when the role exists on paper but is undefined in practice, and
%      it justifies why this project defined the PO's remit explicitly rather
%      than just naming the role.
%   !! This source was originally planned for chapter 5 only. It was moved here
%      because chapter 5 admits no new sources - a source used only in the
%      conclusion is a new source in the conclusion by construction. Using it
%      here makes the chapter-5 back-reference legitimate.
%   !! EVIDENTIAL STATUS: the p. 15 statements are validated learning outcomes;
%      the p. 9 statements are pre-analysis interview perceptions. Use p. 15.
%      Never write "the study found that..." from a p. 9 quote.


% =============================================================================
% 3.2 THE INITIAL PRODUCT BACKLOG AND ITS PRIORITISATION
% -----------------------------------------------------------------------------
% LITERATURE: MIXED  ->  cite the METHOD, never the DECISIONS
%   The rule for this subsection: every concept that has a name gets a source;
%   every judgement about the recipe app is your own and stays uncited.
%   Cite here:
%     - what a product backlog is and what "ordered" means (Scrum Guide)
%     - MoSCoW - it comes from DSDM; find the original or a solid secondary
%       source rather than a blog. The Agile Practice Guide and the SBOK also
%       describe prioritisation techniques
%     - value-based prioritisation, if that route is chosen instead
%     - WSJF, when rejecting it: name what it is with a source, THEN reject it.
%       Rejecting an uncited method looks like ignorance; rejecting a cited one
%       looks like judgement, and that is the Transfer point
%     - non-functional requirements as backlog items, if the optional NFR items
%       are added
%   Research, best case: one academic paper on requirements prioritisation or
%   backlog management in agile projects - it lifts the chapter above textbook
%   level and is exactly the "own research" he expects. The "backlog__*" PDFs
%   in literature/ are that kind of source.
%   NOT cited: the ten priorities themselves, the justification column, the
%   dependencies between your stories.
%
% WHAT YOU WANT TO SAY HERE, IN SHORT
%   - "The items came from user conversations; the initial backlog is exactly
%      the ten stories from section 2.3." [own, one sentence]
%   - "A product backlog is an ordered list of everything known to be needed."
%      [cite]
%   - "MoSCoW was chosen because it is simple and communicates priority to
%      non-technical stakeholders." [cite the method, own the reason]
%   - "WSJF was considered and rejected, because a ten-item backlog does not
%      justify a scoring formula whose inputs are themselves guesses."
%      [cite the method, own the rejection]
%   - "Priority follows business value - not what is technically interesting -
%      and must-haves are ordered among themselves; these dependencies force
%      this order." [own]
%   - "Priority is independent of how well an item is already understood."
%      [own, sets up 3.3]
% -----------------------------------------------------------------------------
% WHERE THE ITEMS COME FROM: one sentence (user conversations, market
% observation). His section plan foresees "how you gather requirements"
% [G-Dec25], but a long origin narrative wastes the page budget.
%
% THE BACKLOG IS EXACTLY THE TEN STORIES. Answering a student who asked how
% many items to include: "There's only one initial backlog that you need. Just
% take the ten user stories that you've come up with and use those as the basis
% for your initial backlog" [G-Jun26]. Asked about epics or "removing noise",
% he described only: build the backlog from the ten stories in priority order
% using a technique such as MoSCoW. => No epic hierarchy, no invented extra
% items. Big page saver.
%   OPTIONAL EXTRA, cheap Transfer point [G-Feb25]: the backlog also holds
%   non-functional requirements "which maybe the engineers will come up with" -
%   one or two NFR items (e.g. scalability, data protection) is a legitimate
%   addition.
%
% YOU NEED NOT WALK ALL TEN THROUGH THE PROSE: "describing a few user stories
% to illustrate how you create the backlog is fine" [T-25.02.2025]. List all
% ten in the table, discuss three or four in the text.
%
% TABLE PLACEMENT: the prioritisation table may go in the appendix
% [T-20.05.2026] - but appendix content counts toward the page limit anyway,
% so it stays in the body.
%
% PRIORITISATION METHOD - ANY STRUCTURED METHOD IS ACCEPTED. Asked "Is it
% possible that we use value benefits approach instead of MoSCoW?" he answered
% "Yeah, you can", adding "any other structured prioritisation approach. The
% easiest one is MoSCoW. That's why I put this down" [G-Jan26]. MoSCoW is a
% convenience default, not a preference. It is separately sanctioned in
% writing: "you can introduce MoSCoW, and indicate how you use it to prioritize
% your backlog" [T-25.02.2025].
%   - MoSCoW's W: he has called it "would have" [G-Mar25], said the distinction
%     "doesn't really matter" [G-Feb25], and glossed it as won't-have
%     [G-Jun26]. Do NOT write a paragraph about it.
%   - Useful justification wording for low priority [G-Jun26]: not in this
%     release because "they're not necessary... too expensive, or too
%     ambitious". His own recipe-app examples of near-zero priority: an East
%     Timor recipes section [G-Sep25] and multi-currency support [G-Jan25].
%   - VALUE MEANS BUSINESS VALUE, not what the engineers find interesting -
%     "a nice little AI widget that the tech guys want to include. Well, if it
%     doesn't offer business value... then it does not get to the top"
%     [G-Jan25].
%   - Prioritise WITHIN the must-haves as well [G-Nov25].
%   - Mention dependencies and constraints as an ordering factor [S4].
%   - Inventing own category criteria (e.g. basic features, personalisation,
%     engagement) was asked in the channel and never answered - neither
%     confirmed nor forbidden. Safe if the criteria are justified.
%   - WSJF: mention it ONLY as a considered-and-rejected alternative, using his
%     own reasoning - "I inherited this slide... I've never seen the weighted
%     shortest job first equation used... it looks too scientific" [G-Feb25];
%     "it's pseudo scientific, so I won't even attempt to explain it and you
%     don't even need to consider it" [G-Jan25]. A rejected alternative with a
%     stated reason is exactly what the Transfer criterion rewards.
%   - PRIORITY AND READINESS ARE INDEPENDENT: "you can have a requirement that
%     is clearly understood, that is estimated, but that is given a really low
%     priority" [G-Jan25]. Useful when annotating the table, and it sets up 3.3.
% =============================================================================
% #############################################################################
% WRITING PLAN - 3.2, PARAGRAPH BY PARAGRAPH
% Budget: 0.80 pages, of which the table takes 0.17 => ~440 WORDS of prose
% across 6 paragraphs => ~70 WORDS EACH. MIXED - about 10 citations.
% THE RULE FOR THIS SUBSECTION: cite the METHOD, own the DECISION. Every concept
% with a name gets a page-numbered source; every judgement about the recipe app
% stays uncited. Citing your own priorities would be wrong, not diligent.
% #############################################################################
\subsection{Das initiale Product Backlog und seine Priorisierung}
\label{subsec:backlog}

% --- P1: provenance (ONE sentence) and what a backlog is ------------------
% TWO CITATIONS.
%
% OWN VOICE, one sentence only: the items came from the user conversations
% named in chapter 1. The tutor's section plan foresees "how you gather
% requirements", but a long origin narrative wastes the page budget.
%
%   \parencite[p.~10]{scrumguide2020}
%   "The Product Backlog is an emergent, ordered list of what is needed to
%    improve the product."
%   -> "Emergent" carries the never-complete argument; "ordered" carries the
%      single-sequence argument. Both are needed later, so quote the whole
%      sentence.
%
%   \parencite[p.~167]{cohn2004}
%   "When a project is initiated there is no comprehensive effort to write down
%    all foreseeable features. Typically, the product owner and team write down
%    everything obvious, which is almost always more than enough for a first
%    sprint."
%   -> Directly licenses a ten-story starting backlog. This is the sentence that
%      pre-empts "why only ten items?".
Die Einträge stammen aus den in Kapitel~1 beschriebenen Nutzergesprächen. Ein
Product Backlog ist \enquote{an emergent, ordered list of what is needed to improve the
product} \parencite[p.~10]{scrumguide2020}. Emergent bedeutet, dass es nie
fertig ist, und ordered bedeutet eine einzige Reihenfolge statt einer Menge von
Kategorien, weshalb die zehn Cookbox-Einträge einen Rang tragen und kein Thema.
Klein anzufangen ist legitim, denn zu Beginn gibt es
\enquote{no comprehensive effort to write down all foreseeable features}, und was
offensichtlich ist, ist \enquote{almost always more than enough for a first sprint}
\parencite[p.~167]{cohn2004}.

Das initiale Backlog von Cookbox besteht damit aus genau den zehn Stories aus
Tabelle~\ref{tab:user-stories} und aus nichts sonst. \textcite[p.~173]{cohn2004} hält
fest, dass \enquote{Scrum has no prescribed, or even recommended, approach to initially
stocking the product backlog}, schließt die Lücke, indem er zuerst die
Nutzerrollen bestimmt --- das Vorgehen aus Abschnitt~\ref{subsec:stories} ---, und
beschränkt das Backlog \enquote{to only user stories}, statt darin auch Fehler und
Untersuchungen zu führen. Diese Einschränkung ist Cohns eigene Abwandlung des reinen Scrum und wird
hier bewusst übernommen.

Als Ordnungsschema diente MoSCoW, die \enquote{must have, should have, could have,
and won't have} method \parencite[PMBOK Guide, Section~4.4.4]{pmbok2021}.
\textcite[p.~98]{cohn2004} übernimmt die MoSCoW-Regeln aus DSDM und füllt die
Bänder inhaltlich: \enquote{Could-have features are ones that can be left out of the
release if time runs out}. US-08 und US-09 stehen in diesem Band und bilden damit
den Puffer dieses Releases gegen einen festen Termin. Vier Kategorien lassen sich
außerdem einem nichttechnischen Stakeholder in einem Satz erklären. Die Schwäche
des Verfahrens wird eingeräumt statt verschwiegen: Eine Übersicht über fünfzehn
Priorisierungsverfahren stellt fest, dass \enquote{most of
the methods are informal and subjective} \parencite[p.~53]{racheva2008}.%
\footnote{Die für diesen Bericht verfügbare Fassung von \textcite{racheva2008}
ist das Autorenmanuskript der Camera-ready-Version, das keine gedruckten
Seitenzahlen trägt. Die hier angegebenen Seitenzahlen sind aus dem
veröffentlichten Bereich 49--58 und den neun fortlaufenden Seiten des
Manuskripts abgeleitet; sie sind als Näherung zu lesen.}

Als Alternative wurde eine Bewertungsformel geprüft. Der Agile Practice Guide
beschreibt eine Rangbildung \enquote{with value including cost of delay divided by duration
(CD3) and other value models} \parencite[p.~59]{pmi2017}, also die Rechnung hinter
Weighted Shortest Job First. Verworfen wurde sie wegen mangelnder Passung, nicht
aus grundsätzlichen Erwägungen. Dieselbe Übersicht
nennt vier Kriterien für die Wahl eines Verfahrens, darunter die Zahl der
Einträge und die verfügbaren Informationsquellen \parencite[p.~56]{racheva2008};
ein Backlog aus zehn Einträgen und eine fehlende gemessene Velocity scheitern an
beiden.

Dieselben zehn Stories erscheinen in Tabelle~\ref{tab:initial-backlog} erneut,
diesmal in Prioritätsreihenfolge und jeweils mit Band und Begründung.


% --- P2: the backlog IS the ten stories, and that is a modification -------
% TWO CITATIONS. This paragraph carries a Transfer point - do not skip it.
%
%   \parencite[p.~173]{cohn2004}
%   "Scrum has no prescribed, or even recommended, approach to initially
%    stocking the product backlog... However, I've found that by first
%    identifying user roles and then focusing on the stories for each user role,
%    Scrum is combined with a powerful requirements identification technique."
%   -> The strongest procedural quote available: it names the exact method used
%      in chapter 2 and ties chapter 2 to chapter 3.
%
%   \parencite[p.~173]{cohn2004}
%   "Rather than allowing product backlog items to describe new features, issues
%    to investigate, defects to be fixed, and so on, the product backlog is
%    constrained to only user stories."
%   -> ** THE TRANSFER POINT **: Cohn presents this as HIS OWN modification to
%      plain Scrum, not as Scrum. Say so explicitly - the report adopts a
%      documented adaptation and names it as such. That is precisely what the
%      15 % Transfer criterion asks for.
%
% OPTIONAL, if a third citation is affordable:
%   \parencite[p.~70]{schwaber2004}
%   "The rest could emerge as we proceeded; we had enough to start with."
%   -> Quote this CLAUSE ONLY, not the flip-chart sentence around it. It
%      justifies a proportionate, deliberately non-exhaustive sprint zero.
%
% OWN VOICE: state that the initial backlog is exactly the ten stories of
% Table~\ref{tab:user-stories} - no epics, no invented extra items.
% OPTIONAL, uncited: one or two non-functional items (e.g. search results under
% one second, offline access to saved recipes). Keep them UNCITED - the sources
% that would support them were dropped, and Cohn actually argues the OPPOSITE
% (he wants constraints on separate cards, not as backlog items).


% --- P3: the prioritisation method ----------------------------------------
% THREE CITATIONS.
%
%   \parencite[PMBOK Guide, Section~4.4.4]{pmbok2021}
%   "Prioritization schema are methods used to prioritize portfolio, program, or
%    project components, as well as requirements, risks, features, or other
%    product information. Examples include a multicriteria weighted analysis and
%    the MoSCoW (must have, should have, could have, and won't have) method."
%   -> ** THE MoSCoW CITATION. ** All four categories, from a PMI standard. This
%      is why the vendor-published SBOK is not needed and why Racheva's version
%      is not used (hers traces to an Oracle white paper).
%
%   \parencite[p.~98]{cohn2004}
%   The substance of the four bands: must-haves are fundamental; should-haves
%   are important but have a short-term workaround; "Could-have features are ones
%   that can be left out of the release if time runs out"; won't-haves are
%   desired but deferred to a later release.
%   !! ATTRIBUTE IT PROPERLY: Cohn borrows MoSCoW from DSDM and footnotes
%      Stapleton (2003). Phrase it "Cohn adopts DSDM's MoSCoW rules".
%   -> Note the "could" band explicitly - "left out if time runs out" is the
%      direct link to the inverted iron triangle in chapter 4. Plant it here.
%
%   \parencite[p.~53]{racheva2008}
%   "The analysis of the methods in the literature indicates that most of the
%    methods are informal and subjective."
%   OPTIONAL companion, the paper's own topic sentence, if 3.2 P3 has room:
%   \parencite[p.~52]{racheva2008} "While in a traditional project the
%   prioritization is usually performed once and before the implementation
%   phase, in agile context it is an ongoing process, performed in the beginning
%   of each iteration, or even during the iteration".
%   -> Lets the report CONCEDE MoSCoW's subjectivity in a cited sentence rather
%      than pretending to rigour. That concession is the honest half of the
%      WSJF rejection in P4.
%   !! Introduce Racheva et al. as a review of fifteen prioritisation methods,
%      never as an empirical study - the paper is explicit work-in-progress.
%
% OWN VOICE: why MoSCoW for THIS project - it is simple, it communicates
% priority to non-technical stakeholders, and the labels are a discussion device
% rather than a score.


% --- P4: WSJF, examined and rejected --------------------------------------
% ONE CITATION. This paragraph is short and scores well above its length.
% Rejecting an UNCITED method looks like ignorance; rejecting a CITED one looks
% like judgement, and that is the Transfer point.
%
%   \parencite[p.~59]{pmi2017}
%   "Rank with value including cost of delay divided by duration (CD3) and other
%    value models"
%   -> CD3 IS the cost-of-delay-divided-by-duration formula behind WSJF. This
%      names the method from a task-mandated source - far safer than citing a
%      SAFe web page for a method you then dismiss.
%
%   \parencite[p.~56]{racheva2008}
%   The four criteria for choosing a prioritisation approach: number of items to
%   be prioritised, number of stakeholders involved, level of requirements
%   volatility, sources of information available.
%   -> ** REST THE REJECTION ON THESE **: a ten-item backlog with one PO and high
%      volatility does not justify a scoring formula whose inputs are themselves
%      guesses. The rejection then follows from stated criteria instead of taste.
%
% !! DO NOT REJECT WSJF ON "value cannot be quantified". Racheva et al.'s own
%    Solution 4 proposes a business-value function; an examiner reading the same
%    paper would find the authors doing what the report says cannot be done.
%    Reject on FIT TO A TEN-ITEM BACKLOG, not on principle.


% --- P5: sentence pointing at the table, then the table -------------------
% NO CITATION.
% One sentence: the ten stories of Table~\ref{tab:user-stories} reappear here in
% priority order as Table~\ref{tab:initial-backlog}. Say out loud that these are
% the SAME ten items - that is the visible difference between question 1 and
% question 2.

% [!ht] rather than [htbp]: the table must stay in the running text of 3.2.
\begin{table}[H]
  \centering
  \caption{Initiales Product Backlog der Rezept-App, priorisiert mit MoSCoW}
  \label{tab:initial-backlog}
  \tablefontsize
  \begin{tabular}{@{}p{1.2cm}p{2.6cm}p{10.5cm}@{}}
    \toprule
    \textbf{ID} & \textbf{Priorität} & \textbf{Begründung} \\
    \midrule
    % ALL TEN stories from Table~\ref{tab:user-stories} and nothing else - the
    % initial backlog is exactly those ten items. Rows are in priority order,
    % and the must-haves are ordered among themselves.
    % Distribution is 4/3/2/1, not the 3/3/3/1 of the original placeholder.
    % Full rationale: PROJECT_FACTS.md, section 5.
    %
    % THREE ROWS DO DOUBLE DUTY - do not reword them without checking 3.3 and
    % chapter 4:
    %   US-07 plants the cost-feedback trigger (3.3 P5)
    %   US-08 plants the legal-obligation trigger (3.3 P5)
    %   US-09 plants the feature that gets dropped in chapter 4, axis 5
    %
    % KEEP EACH JUSTIFICATION TO ONE LINE. This table shares 0.80 pages with six
    % paragraphs of prose, and tables count toward the page limit.
    US-01 & Must have   & Fundament: Sechs der zehn Stories lesen, verändern oder
                          löschen ein Rezept, das diese Story anlegt. \\
    US-02 & Must have   & Das Wiederfinden, nicht das Speichern, ist das Problem,
                          für das es das Produkt gibt. \\
    US-03 & Must have   & Der Einsatzort ist die Küche; ohne brauchbare
                          Kochansicht wird die App einmal probiert und
                          aufgegeben. \\
    US-04 & Must have   & Eine verlorene, von Hand eingegebene Sammlung wird
                          nicht erneut eingegeben; die Bindung hängt daran. \\
    US-05 & Should have & Wertvoll, aber die Wochenplanung ist Komfort und nicht
                          der Kern des Produkts. \\
    US-06 & Should have & Wertvollster Eintrag nach den Must-haves, aber durch
                          US-05 und US-01 blockiert. \\
    US-07 & Should have & Häufig gewünscht; innerhalb des Bandes zurückgestellt,
                          nachdem die Entwickler die Einheitenumrechnung deutlich
                          höher schätzten als zunächst angenommen. \\
    US-08 & Could have  & Zunächst hier eingeordnet; während des Refinements zum
                          Must-have hochgestuft
                          (Abschnitt~\ref{subsec:refinement}). \\
    US-09 & Could have  & Moderation ist bei den Startvolumina manuell
                          bezahlbar; das Werkzeug lohnt sich erst bei Skalierung. \\
    US-10 & Won't have  & Wiederholt gewünscht, aber das Parsen beliebiger
                          Rezeptseiten ist für ein erstes Release
                          unverhältnismäßig teuer. \\
    \bottomrule
  \end{tabular}
\end{table}

% --- P6: interpreting the ordering ----------------------------------------
% THREE CITATIONS. The priorities themselves and the justification column stay
% UNCITED - they are your judgements.
%
% (a) business value, not technical interest:
%   \parencite[p.~41]{pmi2017} "Product owners rank the work based on its
%   business value." - if already used in 3.1, do not repeat it; instead make
%   the point in your own words and name a feature the engineers would have
%   found interesting but that does not reach the top.
%
% (b) dependencies as an ordering factor - TWO citations, one rule and one
%     concrete instance:
%   \parencite[p.~87]{schwaber2004}
%   "Dependencies between requirements are resolved by placing dependent
%    requirements at a position in the list lower than the requirements on which
%    they depend."
%   -> The operational ordering RULE. Trim to this second sentence.
%   \parencite[p.~389]{lucassen2016}
%   "For example, US14 is dependent on US13, because it is impossible to view a
%    person's profile without first laying the foundation for creating a person."
%   -> Transfers directly and concretely: no story about scaling or sharing a
%      recipe can be built before a recipe can be created. Name the two real
%      story IDs from your own table.
%   !! DO NOT use Cohn (2004, p. 17) here. His sentence - "care should be taken
%      to AVOID introducing dependencies" - is the INVEST "Independent"
%      attribute and argues for eliminating dependencies, not for ordering by
%      them. Citing it here would overstate the source.
%
% (c) why deferred items are deferred rather than dropped: own voice. Useful
%     tutor wording to paraphrase - not in this release because they are not
%     necessary, too expensive, or too ambitious.
%
% (d) prioritise WITHIN the must-haves: own voice, one clause.
%
% (e) ** the sentence that sets up 3.3 **: priority says nothing about how well
%     an item is already understood. Own voice, half a sentence, no citation -
%     the empirical backing for it lives in 3.3.
Zwei Überlegungen haben diese Reihenfolge geprägt. Wert bedeutet
Geschäftsnutzen und nicht technisches Interesse: US-10 steht an letzter Stelle,
weil das Parsen beliebiger Rezeptseiten der technisch reizvollste und zugleich
am wenigsten bezahlbare Eintrag der Menge ist. Abhängigkeiten überwiegen dann den
Wert. Abhängige Anforderungen gehören \enquote{at a
position in the list lower than the requirements on which they depend}
\parencite[p.~87]{schwaber2004}. \textcite[p.~389]{lucassen2016} nähern sich
derselben Lage von der anderen Seite: Stories sollen \enquote{schedulable and
implementable in any order} sein, und eine Abhängigkeit, die sich nicht
wegentwerfen lässt --- etwa ein Profil, das sich nicht ansehen lässt
\enquote{without first laying the foundation for creating a person} ---, soll
zumindest ausdrücklich benannt werden. US-01 ist genau
dieser unvermeidbare Fall. Die Story steht aus diesem Grund an der Spitze und
nicht wegen ihres Werts, den US-02 wohl eher trägt: US-02, US-03, US-04, US-06,
US-07 und US-08 lesen, verändern oder löschen alle ein Rezept, das US-01 anlegt.
Auch die Must-haves sind untereinander geordnet und nicht als Block belassen.
US-10 ist zurückgestellt, nicht aufgegeben. Die Priorität sagt zudem nichts
darüber aus, wie gut ein Eintrag bereits verstanden ist, was
Abschnitt~\ref{subsec:refinement} aufgreift.


% =============================================================================
% 3.3 BACKLOG REFINEMENT WITH THE SCRUM TEAM
% This subsection carries the literal wording of the task question and is,
% together with 3.4, the emphasis of Task 2. Refinement comes first.
% -----------------------------------------------------------------------------
% LITERATURE: MIXED, leaning literature  ->  define with sources, illustrate
%                                            with your own project
%   This subsection answers the sentence that makes Task 2 different from
%   Task 1, so the underlying concept must be visibly researched. He said it
%   outright: "you need to research what these are and describe what these
%   are" [G-Jun26].
%   Cite here:
%     - refinement as an ongoing activity, and what it does to backlog items
%       (Scrum Guide; note the 2020 edition treats it as an activity, not an
%       event - a precise, page-numbered citation is worth having)
%     - progressive elaboration as a concept (PMBOK 7 uses the term; it lets
%       you connect the agile practice to a classical PM model, which is
%       exactly what "Transfer" rewards)
%     - the claim that there is no fixed cadence - if a source can be found,
%       cite it; otherwise present it as a reasoned choice, not as a fact
%     - Kelly (2019) for the PO's side of the negotiation
%   Research, best case: one empirical paper on how backlogs actually evolve,
%   or on requirements churn in agile projects, to back the statement that a
%   backlog changes continuously. Avoid the DoR/DoD literature - Gannon never
%   uses the term and it eats space.
%   NOT cited: the nutritional-value example as applied to your app, your
%   chosen cadence, your triggers.
%
% WHAT YOU WANT TO SAY HERE, IN SHORT
%   - "Refinement, also called grooming, is the continuous activity of adding,
%      removing, splitting, re-estimating and re-ordering backlog items."
%      [cite]
%   - "It is joint work: the PO knows the business priority, the developers
%      know their capacity, and the two are negotiated." [own, cite Kelly]
%   - "Worked example: a coarse request enters the backlog as a block, and is
%      progressively broken down until it is small, clear and estimable enough
%      to be pulled into a sprint." [own example, cite the concept]
%   - "Readiness is not a document - it means small enough, clear enough,
%      estimated enough." [own]
%   - "These events trigger refinement in this project: the sprint review, a
%      new legal requirement, user feedback that demotes a must-have." [own]
%   - "The cadence chosen here is X, because the team is new to the domain;
%      there is no universal rule." [own, cite if possible]
% -----------------------------------------------------------------------------
% WHAT REFINEMENT IS - his own definitions, usable almost verbatim as a basis:
%   "repeated conversations between the product owner, the development team,
%    and users to ask: is this the right priority? ... Are we happy with the
%    estimates?" [G-Mar25]
%   "the continuous changing of the backlog to account for new changes, new
%    business requirements, or different prioritisation" [G-Nov25]
%   compact version: "introducing new requirements, picking out redundant
%    requirements, and reprioritization... and identifying what needs to go
%    into the sprint" [G-Feb25]
% Say that refinement is also called GROOMING [G-Nov25]. Stress that it is an
% ongoing activity, not a single scheduled meeting.
%
% WHY IT IS JOINT - the negotiation logic, and the reason the task says "with
% the Scrum team" [G-Feb25]: "the product owner knows what is of priority to
% the business... The technologists know what their capacity is... there's a
% kind of negotiation that goes on."
%
% PROGRESSIVE ELABORATION - his best recipe-app example, worth reproducing in
% own words [G-Jan25]: a user asks for "categorising nutritional value of
% food"; that "might require ten different user stories"; it enters the backlog
% as a high-level block, is refined over time and the resulting items "move up
% the stack... until they get to a point where they're clearly defined and
% estimated".
%
% READINESS TEST IN PLAIN WORDS, not DoR jargon - [G] never uses DoR/DoD;
% they are [S]-only. Keep any mention very tight. His plain formulation
% [G-Feb25]: small enough, clearly defined enough, estimated enough.
%
% OTHER TRIGGERS TO NAME:
%   - the SPRINT REVIEW as the recurring trigger - "one key trigger for
%     refinement... at the end of every sprint" [G-Sep25]
%   - a new LAW mid-project pushing an item to the top [G-Mar25]
%   - a must-have DOWNGRADED after user feedback [G-Nov25]
%   - stories added and removed [G-Mar25]
%   Pick recipe-app instances for each - generic trigger lists score nothing.
%
% CADENCE - present it as a deliberate decision and state that there is no
% fixed rule ("no consensus how long the refinement should be" [S3]). Options
% [S3, S4]: just-in-time; a time-boxed session of about an hour midway through
% the sprint; several sessions when the team is new to the domain - the
% defensible pick for a fresh recipe app. Justify the choice.
%
% Priority and maturity are independent of each other - refer back to 3.2 in
% half a sentence, do not restate the point.
% =============================================================================
% #############################################################################
% WRITING PLAN - 3.3, PARAGRAPH BY PARAGRAPH
% Budget: 0.70 pages = ~490 WORDS across 7 paragraphs => ~70 WORDS EACH.
% MIXED, leaning cited - about 11 citations.
%
% ** THIS SUBSECTION CARRIES THE LITERAL WORDING OF THE TASK QUESTION. ** It is
% what makes this Task 2 rather than Task 1, and the tutor said outright: "you
% need to research what these are and describe what these are."
%
% Sources that own it: scrumguide2020 (the definition), pmi2017 pp. 52-53 (the
% most concrete refinement passage in the whole corpus - cadence AND a time
% budget), heikkila2017 (one empirical figure), schoen2017, racheva2008.
%
% !! DO NOT reach for schwaber2004 or cohn2004 for the DEFINITION - neither book
%    contains "backlog refinement" at all. Cohn's word "refinement" always means
%    successive refinement of a system or a role list. Use Cohn for the
%    SUBSTANCE (sizing, splitting, deferring detail) only.
% #############################################################################
\subsection{Backlog Refinement mit dem Scrum Team}
\label{subsec:refinement}

% --- P1: what refinement is ------------------------------------------------
% THREE CITATIONS plus one footnote.
%
%   \parencite[p.~10]{scrumguide2020}
%   "Product Backlog refinement is the act of breaking down and further defining
%    Product Backlog items into smaller more precise items. This is an ongoing
%    activity to add details, such as a description, order, and size."
%   -> THE definition sentence. Note the Guide's word is "ongoing" - do not
%      write "continuous" inside the quotation marks (0 hits in the Guide).
%
%   \parencite[p.~150]{pmi2017}
%   "Backlog Refinement. The progressive elaboration of project requirements
%    and/or the ongoing activity in which the team collaboratively reviews,
%    updates, and writes requirements to satisfy the need of the customer
%    request."
%   -> Every clause is usable: progressive elaboration, ongoing ACTIVITY,
%      collaboratively, and three distinct actions. This is also the citation
%      that answers "with the Scrum team" head-on.
%
%   \parencite[p.~7]{scrumguide2020}
%   "Events are used in Scrum to create regularity and to minimize the need for
%    meetings not defined in Scrum."
%   -> ** POSITIVE evidence that refinement is not a ceremony **, rather than an
%      argument from silence. The Guide defines no refinement meeting, timebox
%      or attendee list, and this sentence explains why. It also licenses
%      choosing your own cadence, which is the Transfer move in P6.
%
%   FOOTNOTE on terminology - cheap Documentation point:
%   \parencite[ch.~8, pp.~57--62]{kelly2019}
%   "In some circles this exercise has been called grooming. However, in
%    England, this term has taken on negative connotations with the general
%    public in recent years and is better avoided."
%   -> Use this to explain in a footnote why the report says "refinement" and
%      not "grooming".
%   !! DO NOT attribute the refinement = grooming equivalence to the Scrum
%      Guide - "grooming" occurs 0 times in it.
Product Backlog Refinement ist \enquote{the act of breaking down and further
defining Product Backlog items into smaller more precise items}
\parencite[p.~10]{scrumguide2020}.\footnote{\textcite[ch.~8,
pp.~57--62]{kelly2019} hält fest, dass die ältere Bezeichnung Grooming in England
\enquote{has taken on negative connotations with the general public} und
\enquote{is better avoided}.} Der Guide bezeichnet es als fortlaufend und sieht
dafür kein Event vor; Events existieren \enquote{to create regularity and to
minimize the need for meetings not defined in Scrum} \parencite[p.~7]{scrumguide2020}. Die Taktung bei Cookbox ist
deshalb eine Entscheidung und keine übernommene Regel.

Das Verfeinern der zehn Stories ist gemeinsame Arbeit. In agilen Vorhaben gilt:
\enquote{the product owners create the backlog for and with the team}
\parencite[p.~41]{pmi2017}. Die Entwickler, die die Arbeit ausführen werden,
\enquote{are responsible for the sizing}, und der Product Owner hilft ihnen,
\enquote{understand and select trade-offs}
\parencite[p.~11]{scrumguide2020}. Die zeitliche Abhängigkeit aus Kundensicht
\enquote{could be different from chronological dependency from developers'
perspective} \parencite[p.~55]{racheva2008}. US-06 ist hier der Fall: Die Nutzer
wollten die Einkaufsliste früh, die Entwickler ordneten sie zunächst hinter
US-05 und US-01 ein.

Refinement macht grobe Wünsche bearbeitbar, was PMI progressive elaboration
nennt: \enquote{the iterative process of increasing the level of detail}, sobald
genauere Schätzungen verfügbar werden \parencite[p.~153]{pmi2017}. Einträge nahe
der Spitze werden so geschnitten, dass sie in eine Iteration passen,
\enquote{while more distant stories
could be much larger and less precise} \parencite[p.~78]{cohn2004}. Eine
Teilnehmerin bat um Unterstützung dabei, das aufzubrauchen, was in ihrem
Kühlschrank bald verdirbt. Der Wunsch lässt sich nicht schätzen und enthält
plausibel acht bis zehn Stories, weshalb er als einzelner Block unterhalb der
zehn ins Backlog aufgenommen wurde.

Die Reife war ein einfacher Test und kein Dokument. Eine Sitzung verfeinert so
viele Stories, dass das Team versteht, \enquote{what the stories are and how large the stories
are in relation to each other} \parencite[p.~52]{pmi2017}, und hält sie
\enquote{small enough so the team can produce a steady flow of completed work}
\parencite[p.~53]{pmi2017}. Klar genug, geschätzt genug, klein genug: Der
Kühlschrank-Block scheiterte an allen dreien, US-01 und US-02 bestanden.

Vier Anlässe haben das Backlog verändert. Eine datenschutzrechtliche Prüfung
stellte fest, dass die Kontolöschung eine gesetzliche Pflicht ist und kein
Komfortmerkmal, weshalb US-08 von Could have auf Must have stieg, oberhalb von
drei höher eingestuften Einträgen. Die Schätzung von US-07 zeigte, dass die
Einheitenumrechnung über Gramm, Milliliter, Tassen und ganze Eier hinweg ihre
einzeilige Beschreibung weit übersteigt, weshalb sie an das Ende des
Should-have-Bandes rückte --- eine Kostenrückmeldung, die eine Priorität
verändert, wie \textcite[p.~100]{cohn2004} sie beschreibt. Nachfolgende Gespräche zeigten, dass die
meisten Teilnehmer ohnehin auf Papier planen, weshalb US-06 so umformuliert
wurde, dass die Liste aus ausgewählten Rezepten entsteht; damit rutschte US-05
in der Reihenfolge darunter und die Abhängigkeit von US-05 entfiel. Der
Kühlschrank-Block wurde über zwei Sitzungen aufgeteilt.

Die Taktung wurde entschieden und nicht nachgeschlagen. PMI berichtet von
\enquote{no consensus on how
long the refinement should be}, hält fest, dass viele iterationsbasierte Teams
\enquote{a timeboxed 1-hour discussion midway through a 2-week iteration} nutzen,
und billigt mehrere Gespräche für Teams, die \enquote{new to the product, the product
area, or the problem domain} sind \parencite[p.~52]{pmi2017}; die genannte
Obergrenze lautet \enquote{not more than 1 hour per week} \parencite[p.~53]{pmi2017}. Dieses Team
ist mit allen dreien neu, weshalb zwei einstündige Sitzungen je zweiwöchigem
Sprint angesetzt wurden: das Doppelte der gängigen Praxis und genau auf der
Obergrenze.

Priorität und Reife sind verschiedene Urteile. In einer Expertenbefragung wird
\enquote{discussing the maturity level of a requirement with the team} als
eigenständige Technik vorgeschlagen, neben dem fortlaufenden Priorisieren des
Backlogs \parencite[p.~46]{schoen2017}, was auf zwei unabhängige Dimensionen
hindeutet. US-08 zeigt, warum: Die rechtliche Prüfung hob die Priorität, aber
der Eintrag war deshalb nicht besser verstanden und musste vor der Auswahl
weiterhin verfeinert werden. Die in Abschnitt~\ref{subsec:backlog} festgelegte
Reihenfolge ist damit noch kein Plan.


% --- P2: why refinement is JOINT -------------------------------------------
% THREE CITATIONS. This paragraph is the direct answer to "refined WITH the
% Scrum team" and must not be thin.
%
%   \parencite[p.~41]{pmi2017}
%   "In agile, the product owners create the backlog for and with the team."
%   -> Quote it verbatim; the two prepositions ARE the point.
%
%   \parencite[p.~11]{scrumguide2020}
%   "The Developers who will be doing the work are responsible for the sizing.
%    The Product Owner may influence the Developers by helping them understand
%    and select trade-offs."
%   -> This IS the priority-versus-capacity negotiation, stated as a division of
%      labour. The PO knows what the business needs; the developers know what
%      they can carry.
%   !! ***** DECIDED: the Scrum Guide p. 11 sizing sentence is spent HERE, not
%      in 3.4. ***** This subsection carries the literal task wording ("refined
%      with the Scrum team") and is the graded difference between Task 1 and
%      Task 2, so the sentence that states the division of labour belongs here.
%      3.4 makes the same point in its own words and does not re-cite it.
%
%   \parencite[p.~55]{racheva2008}
%   "A chronological dependency from clients' point of view could be different
%    from chronological dependency from developers' perspective, as the first one
%    is based on business reasons... while the second one has technical reasons
%    and might require completely different sequence of implementing features."
%   -> ** THE MECHANISM **: the two sides hold two different orderings of the
%      same items, so the ordering cannot be settled by either alone. This is
%      the strongest available answer to "why joint?" and it goes beyond the
%      normative sources, which only say that it IS joint.
%
% NOTE: the Agile Manifesto's principle 4 ("Business people and developers must
% work together daily") was considered here and REJECTED. The first pass called
% it irreplaceable; it is not - the three citations above cover the same ground
% with page numbers, and the Manifesto's three citations are committed to
% chapters 4 and 5.


% --- P3: progressive elaboration, with a worked recipe-app example --------
% TWO CITATIONS for the concept; the EXAMPLE is your own and stays uncited.
%
%   \parencite[p.~153]{pmi2017}
%   "Progressive Elaboration. The iterative process of increasing the level of
%    detail in a project management plan as greater amounts of information and
%    more accurate estimates become available."
%   -> Naming the classical PM term and then showing the agile practice doing it
%      is a Transfer point: it connects the two paradigms rather than treating
%      them as unrelated.
%
%   \parencite[p.~78]{cohn2004}
%   "This means, for example, that stories for the next few iterations would be
%    written at sizes that can be planned into those iterations, while more
%    distant stories could be much larger and less precise."
%   -> The single best quote for "top of the backlog small and clear, bottom
%      coarse".
%
% OWN VOICE - the worked example. Take a coarse request from a real user
% conversation (the tutor's own recipe-app instance is "categorising nutritional
% value of food", which might require ten separate stories). Show it entering
% the backlog as one high-level block, being broken down over successive
% refinement sessions, and the resulting items rising until they are small,
% clear and estimable enough to be pulled into a sprint.
% !! Pick a request that is NOT already one of the ten stories, and make sure it
%    does not collide with the out-of-scope list in 2.1 (nutrition may be
%    excluded there - if so, choose another coarse request, e.g. "help me plan a
%    week of meals").


% --- P4: readiness, in plain words ----------------------------------------
% TWO CITATIONS. Keep DoR jargon out - the tutor never uses the term.
%
%   \parencite[p.~52]{pmi2017}
%   "The purpose of these meetings is to refine enough stories so the team
%    understands what the stories are and how large the stories are in relation
%    to each other."
%   -> Gives "clear enough" and "estimated enough" in PMI's own words.
%
%   \parencite[p.~53]{pmi2017}
%   "...making sure all of the stories are small enough so the team can produce a
%    steady flow of completed work."
%   -> Gives "small enough".
%
%   OPTIONAL exit criterion, if a third is affordable:
%   \parencite[p.~10]{scrumguide2020}
%   "Product Backlog items that can be Done by the Scrum Team within one Sprint
%    are deemed ready for selection in a Sprint Planning event."
%   (Note the capital D on "Done" in the original.)
%
% !! The literature review's first pass proposed PMI's Definition of Ready
%    (2017, p. 151) here. It has been CUT: the outline wants the plain-words
%    formulation, and the three quotes above already deliver small enough /
%    clear enough / estimated enough without importing the acronym.


% --- P5: what triggers refinement ------------------------------------------
% TWO CITATIONS; the recipe-app INSTANCES are your own. A generic trigger list
% scores nothing - each trigger needs a named consequence for a named story.
%
%   \parencite[p.~9]{scrumguide2020}
%   "Based on this information, attendees collaborate on what to do next. The
%    Product Backlog may also be adjusted to meet new opportunities."
%   -> The Sprint Review as the recurring trigger, and the cleanest evidence
%      that planning is rolling rather than fixed.
%
%   \parencite[p.~100]{cohn2004}
%   "It was a high priority to her when she thought it was a few hours; when it
%    was a week, there were many other things she decided she would rather have."
%   -> Cost feedback from the team changes the PO's priority - a concrete,
%      mechanism-level trigger rather than a category label.
%
%   OPTIONAL, for the "requirements are not fixed at the start" claim:
%   \parencite[p.~44]{schoen2017} challenge C4 - continuous requirements
%   management is a challenge because not all requirements are fixed at the
%   beginning and may change over the project. Report the denominator honestly:
%   16 of the 22 experts who answered.
%
% OWN INSTANCES to name (pick two or three, each tied to a story ID):
%   - a new legal obligation arriving mid-project and pushing an item to the top
%   - a must-have demoted after user feedback
%   - a story added, and one removed as redundant


% --- P6: cadence as a deliberate decision ---------------------------------
% THREE CITATIONS. ** THIS IS THE HIGHEST-VALUE TRANSFER MOVE IN CHAPTER 3. **
%
% The argument to build, in this order:
%
% (a) there is no rule:
%   \parencite[p.~52]{pmi2017} "There is no consensus on how long the refinement
%   should be."
%   -> Because PMI says there is no consensus, the report MUST justify its own
%      choice. That is the licence for everything that follows.
%
% (b) the common practice:
%   \parencite[p.~52]{pmi2017} "Many iteration-based agile teams use a timeboxed
%   1-hour discussion midway through a 2-week iteration."
%
% (c) the budget ceiling:
%   \parencite[p.~53]{pmi2017} "Teams often have a goal of spending not more than
%   1 hour per week refining stories for the next batch of work."
%
% (d) the warrant for MORE than one session:
%   \parencite[p.~52]{pmi2017} "Multiple refinement discussions for
%   iteration-based agile teams. Teams can use this when they are new to the
%   product, the product area, or the problem domain."
%
% ** THE MOVE **: (b) and (c) look contradictory but reconcile as common practice
% versus budget ceiling - up to two hours per two-week sprint. A team new to the
% recipe domain qualifies under (d) for multiple sessions, and two one-hour
% sessions per two-week sprint sits EXACTLY on PMI's own ceiling. Argue it that
% way and the cadence is DERIVED, not asserted.
% !! Quote (b) and (c) together or neither - quoting only one of the two looks
%    like a misreading if the examiner opens the other page.
%
% (e) OPTIONAL - the empirical target state, if a fifth citation is affordable:
%   \parencite[p.~2912]{heikkila2017}
%   "The goal is to have at least two sprints worth of work for each team in
%    their team backlog at all times."
%   -> The only OBSERVED target stock of refined work in the corpus; PMI stops
%      at "sufficient items for the next iteration", i.e. one. The mild tension
%      between one and two is itself a Transfer point.
%   !! Say "two sprints' worth of WORK", not "of READY work" - the source does
%      not make the readiness claim.
%   !! Pre-empt the scale objection in one clause, using the paper's own words:
%      the case organisation's process was "inspired by the single-team Scrum
%      process" \parencite[p.~2917]{heikkila2017}. Citing that alongside turns a
%      weakness into evidence of source-critical reading.
%
% (f) OPTIONAL - cadence really is a choice, not a standard:
%   \parencite[p.~46]{schoen2017} practitioners report both "refining and
%   prioritizing continuously the product backlog" and "grooming on demand".
%   -> The empirical counterweight to Heikkila's fixed fortnightly rhythm.


% --- P7: priority and maturity are independent (half a sentence) ----------
% ONE CITATION. Refer back to 3.2 - do not restate the point at length.
%
%   \parencite[p.~46]{schoen2017}
%   "discussing the maturity level of a requirement with the team"
%   -> Practitioners report this as a distinct activity ALONGSIDE continuous
%      prioritisation, which is the empirical backing for treating priority and
%      readiness as independent dimensions - and it shows maturity is negotiated
%      with the team rather than decreed by the PO.
%   !! HEDGE IT CORRECTLY. The paper never asserts that priority and maturity are
%      independent. Write: "practitioners report discussing the maturity level of
%      a requirement with the team as a distinct activity alongside continuous
%      prioritisation (Schoen et al., 2017, p. 46)", and let the independence
%      claim be the report's own inference. Overstating this into "the literature
%      shows priority and readiness are independent" is the one place in this
%      subsection where an attribution error is catchable.


% =============================================================================
% 3.4 SPRINT PLANNING
% -----------------------------------------------------------------------------
% LITERATURE: MIXED  ->  the event is defined by the framework, the sprint is
%                        yours
%   Cite here:
%     - the three sprint planning questions and the composition of the sprint
%       backlog, goal + PBIs + plan (Scrum Guide, with page number - this is
%       the single most quotable passage of the whole report)
%     - the one-day task granularity, if stated
%     - estimation as an activity of the team rather than of a manager - a
%       short citation is enough; do not open an estimation chapter
%   Research, best case, and only if space allows: one paper on why new teams
%   estimate their own capacity poorly early on. That would turn the
%   "deliberately conservative first sprint" from an assertion into an
%   evidence-based decision - a strong Process point. The "estimation__*" PDFs
%   in literature/ are that kind of source. If space is tight, argue it
%   yourself in one sentence and move on.
%   Do NOT: build a planning-poker section with citations. It belongs to
%   Task 3 [G-Jun26].
%
% WHAT YOU WANT TO SAY HERE, IN SHORT
%   - "Sprint planning answers three questions: why is this sprint valuable,
%      what can be delivered, and how will it be built." [cite]
%   - "For this project the sprint goal is X, and these items were selected
%      from the top of the backlog to serve it." [own - name real story IDs]
%   - "The selected items were decomposed into tasks of a day or less; the
%      sprint backlog is the goal, the items and the plan, not just a list."
%      [cite the definition, own the content]
%   - "The team estimated the selected items itself, and the first sprint was
%      deliberately underloaded because a new team cannot yet judge its
%      capacity." [own reasoning, optional citation]
%   - Close the chapter here: the project stops at the end of sprint zero.
% -----------------------------------------------------------------------------
% STRUCTURE IT ON THE THREE SCRUM GUIDE QUESTIONS [S3]:
%   1. Why is this sprint valuable? The PO proposes the value, which becomes
%      the SPRINT GOAL.
%   2. What can be done this sprint? The developers select PBIs from the
%      backlog.
%   3. How will the chosen work get done? Decomposition into tasks of one day
%      or less - "we are not talking about user stories any longer".
% SPRINT BACKLOG = sprint goal + selected PBIs + the plan for delivering them,
% "not just a list" [S3].
%
% ESTIMATION IS LICENSED HERE, briefly: "how does the scrum team work to refine
% and estimate what goes into the sprint backlog?" [G-Jan26]. A few sentences,
% located inside sprint planning.
%
% PLANNING POKER: A PASSING MENTION AT MOST, never a section of its own - the
% ruling hardened over time:
% "you can talk about planning poker if you want, if you've got space"
% [G-Sep25] -> "planning poker is not an element of task one and task two. So
% it's unique to Task three" [G-Jun26]. A passing mention at most.
%   IF it is mentioned, one short CONCRETE paragraph beats a vague long one
%   [SCH: name the number range, the reference stories, and how the team
%   reached agreement]. Material: story points are abstract, relative to a
%   reference story and measure complexity rather than time - "23 does not mean
%   23 days of effort" [G-Mar25]; simultaneous reveal draws on the wisdom of
%   crowds so "it doesn't allow the bias of a very strong voice to dominate"
%   [G-Jan26]; the infinity card means the story is too vague and must be
%   rewritten or split [G-Jan25] - the one clean link back to refinement;
%   critical nuance if depth is wanted: it is "prone to group dynamic
%   processes, such as group think" [S3].
%
% NEW TEAMS ESTIMATE CAPACITY BADLY AT FIRST and improve [G-Jan25] - use this
% to justify a deliberately conservative first sprint. This is a Process point:
% it shows the plan accounts for its own uncertainty.
%
% The report ENDS HERE, at the initial sprint. No sprint 2, no release plan.
% =============================================================================
% #############################################################################
% WRITING PLAN - 3.4, PARAGRAPH BY PARAGRAPH
% Budget: 0.40 pages = ~280 WORDS across 4 paragraphs => ~70 WORDS EACH.
% The shortest subsection in chapter 3. About 6 citations,
% five of them the Scrum Guide, because the event is defined by the framework
% and the sprint is yours.
%
% !! PMI's glossary must NOT be used for the sprint backlog: it calls it "A list
%    of work items" (2017, p. 154), which contradicts the goal + PBIs + plan
%    definition this subsection needs. Scrum Guide only.
% !! PMBOK 7 has no sprint-planning model at all - one synonym sentence and
%    nothing on the three topics. Do not reach for it here.
% #############################################################################
\subsection{Sprint Planning}

% --- P1: the three topics, and the sprint goal ----------------------------
% TWO CITATIONS.
%
%   \parencite[p.~8]{scrumguide2020}
%   "Topic One: Why is this Sprint valuable?" / "Topic Two: What can be Done
%    this Sprint?" / "Topic Three: How will the chosen work get done?"
%   -> Structure the whole subsection on these three. Note the capital D on
%      "Done" in Topic Two if quoting verbatim.
%
%   \parencite[p.~8]{scrumguide2020}
%   "The Product Owner proposes how the product could increase its value and
%    utility in the current Sprint. The whole Scrum Team then collaborates to
%    define a Sprint Goal that communicates why the Sprint is valuable to
%    stakeholders."
%
% OWN VOICE: state the actual sprint goal for this project in one sentence.
Das Sprint Planning eröffnet den Sprint und ist um drei Themen herum angelegt:
warum der Sprint wertvoll ist, was in ihm getan werden kann und wie die
ausgewählte Arbeit erledigt wird \parencite[p.~8]{scrumguide2020}. Der Product
Owner schlägt vor, wie das Produkt seinen Wert steigern könnte, und das gesamte
Scrum Team legt anschließend ein Sprintziel fest, das den Stakeholdern vermittelt,
warum der Sprint wertvoll ist \parencite[p.~8]{scrumguide2020}. Als Ziel des ersten Sprints wurde gesetzt,
dass eine kochende Person ein Rezept in Cookbox ablegen und wiederfinden kann.

Welche Einträge diesem Ziel dienen, entscheidet nicht der Product Owner. Der
Guide ist darin ausdrücklich: \enquote{Through discussion with the Product Owner, the Developers
select items from the Product Backlog to include in the current Sprint}
\parencite[p.~8]{scrumguide2020}. Die vier Entwickler nahmen US-01 und US-02 von
der Spitze der Tabelle~\ref{tab:initial-backlog} und schätzten sie selbst, gemäß
der in Abschnitt~\ref{subsec:refinement} dargelegten Arbeitsteilung. Ein Rezept
einzugeben und es wiederzufinden ist der kürzeste Weg zu einem vorführbaren
Ziel.

Das dritte Thema ist der Plan. Die beiden Einträge wurden in Arbeitspakete von
höchstens einem Tag zerlegt, was der Guide als übliche Praxis und nicht als Regel
darstellt \parencite[p.~8]{scrumguide2020}. Daraus entsteht das Sprint Backlog,
\enquote{composed of the Sprint Goal (why), the set of Product Backlog items
selected for the Sprint (what), as well as an actionable plan for delivering the
Increment (how)} \parencite[p.~11]{scrumguide2020}. Bei Cookbox enthält es das
Ziel, US-01 und US-02 sowie die Aufgaben, in die die Entwickler sie zerlegt
haben: einen Plan und nicht bloß eine Liste.

Zwei der vier Must-haves blieben im zweiwöchigen Sprint außen vor. US-03 und
US-04 rangieren über jedem Should-have, aber das Team hat nie zuvor zusammen
gearbeitet und verfügt über keine Velocity, gegen die es planen könnte, und
Teams, die mit Scrum beginnen, übernehmen sich im ersten Sprint häufig
\parencite[p.~107]{schwaber2004}. Bewusstes Unterladen ist billiger, als das
Ziel zu verfehlen. Hier endet die Planung: Das Projekt hört bei einem geplanten
ersten Sprint auf, ohne dass bereits ein Increment gebaut wäre.


% --- P2: what can be done - the developers select -------------------------
% ONE CITATION, plus one conditional.
%
%   \parencite[p.~8]{scrumguide2020}
%   "Through discussion with the Product Owner, the Developers select items from
%    the Product Backlog to include in the current Sprint. The Scrum Team may
%    refine these items during this process, which increases understanding and
%    confidence."
%   -> Two things at once: the DEVELOPERS select (not the PO), and refinement
%      continues inside sprint planning - which is the clean link back to 3.3.
%   !! DO NOT cite Schwaber (2004, p. 134) here. It says the opposite - that the
%      decision of what constitutes the sprint is the PO's responsibility - and
%      citing both would be a visible self-contradiction.
%
%   ESTIMATION - ***** DECIDED: the citation was spent in 3.3 P2. *****
%   Do NOT cite \parencite[p.~11]{scrumguide2020} again here. Make the point in
%   the report's own words - the four developers sized the two selected items
%   themselves - and refer back to 3.3 in half a clause.
%   -> Note the Guide's word is "sizing"; it never uses "estimate" or
%      "estimation" (0 hits). Do not open an estimation section - effort
%      estimation and planning poker belong to Task 3.
%
% OWN VOICE: name the actual story IDs pulled from the top of
% Table~\ref{tab:initial-backlog}, and say why those serve the sprint goal.


% --- P3: how it gets done, and what the sprint backlog is -----------------
% TWO CITATIONS.
%
%   \parencite[p.~8]{scrumguide2020}
%   "This is often done by decomposing Product Backlog items into smaller work
%    items of one day or less."
%   -> Note "often done", not mandated - do not harden it into a rule.
%
%   \parencite[p.~11]{scrumguide2020}
%   "The Sprint Backlog is composed of the Sprint Goal (why), the set of Product
%    Backlog items selected for the Sprint (what), as well as an actionable plan
%    for delivering the Increment (how)."
%   -> The better of the Guide's two sprint-backlog definitions, because it
%      labels why/what/how explicitly and maps straight onto the three topics.
%      Make the point that the sprint backlog is NOT just a list.
%   !! There is a second, weaker definition on p. 9. Cite ONE; citing both with
%      different pages looks like a slip.


% --- P4: the deliberately conservative first sprint -----------------------
% TWO CITATIONS, and this is where the subsection earns its Process marks: it
% shows the plan accounting for its own uncertainty.
%
%   \parencite[p.~8]{scrumguide2020}
%   "Selecting how much can be completed within a Sprint may be challenging."
%
%   \parencite[p.~107]{schwaber2004}
%   "As happens in most organizations starting to use Scrum, many of Service1st's
%    teams overcommitted themselves for the first Sprint."
%   -> The only page in the corpus stating over-commitment as a GENERAL
%      observation across organisations rather than as one anecdote.
%   !! DO NOT use Schwaber (2004, p. 19) for this claim. Read in full, he
%      minimised first-sprint functionality because of architecture and
%      infrastructure work, not because capacity was unknown - a different
%      argument, and a catchable misreading.
%
%   OPTIONAL, if a third citation is affordable - the only quantified statement
%   available, which turns the decision into an inference from evidence:
%   \parencite[p.~64]{pmi2017}
%   "Teams might discover it can take four to eight iterations to achieve a
%    stable velocity."
%
% OWN VOICE: the first sprint was deliberately underloaded, and by roughly how
% much. Then CLOSE THE CHAPTER - the project stops at the end of sprint zero.
% No sprint 2, no release plan, no burndown.
